\clearpage\SourceLeaf{182}
\begin{center}
{\Large BRIEFWECHSEL.}
\end{center}

\bigskip
\begin{center}
[DIE TRANSZENDENTE TRIGONOMETRIE.]
\end{center}

\begin{center}\rule{2cm}{0.4pt}\end{center}

\begin{center}
[\textsc{Wachter} \emph{an} \textsc{Gauss}. \emph{Danzig, 12. Dezember 1816.}]
\end{center}

\{.\ .\ .\ .\ .\ .\ .\ .\ Also wäre die antieuklidische Geometrie oder Ihre Geometrie wahr.
Ihre \textls{Konstante} bleibt aber unbestimmt. Warum? Das ließe sich vielleicht durch
Folgendes verständlich machen. .\ .\ .\ .\ .\ .\ .

.\ .\ .\ .\ .\ .\ .\ .\ Das bisherige Ergebnis wäre also so auszusprechen.

Die
euklidische Geometrie ist falsch. Dennoch muss die wahre Geometrie mit
demselben elften euklidischen Axiom oder mit dem Postulat von Linien und Flächen
beginnen, die die in jenem Axiom behauptete Eigenschaft haben. An die Stelle der
Geraden und der Ebene sind nur die Großkreise jener mit unendlichem Radius beschriebenen Kugel
und ihre Oberfläche zu setzen. Daraus entsteht zwar die eine
Unbequemlichkeit, dass die Teile dieser Fläche bloß symmetrisch und anders als in
der Ebene nicht kongruent sind oder dass der Radius nach der einen Seite hin
unendlich und nach der anderen imaginär ist. Es ist aber klar, wie diese Unbequemlichkeit durch
viele andere Vorteile aufgewogen wird, die die Konstruktion auf einer Kugelfläche bietet,
\textls{sodass vielleicht auch dann noch}, wenn die
euklidische Geometrie wahr wäre, zwar keine Notwendigkeit mehr bestünde,
die Ebene als eine unendliche Kugelfläche zu betrachten, aber doch noch die
Fruchtbarkeit dieser Sichtweise dafür sprechen könnte.

Als ich aber all dies durchdacht und mich mit dem Ergebnis schon
völlig abgefunden hatte, teils weil ich glaubte, der Grund (la métaphysique)
jener der Geometrie notwendig anhaftenden Unbestimmtheit --- und sogar
